
\section{Introduction}\label{sec:intro}
The market for US-dollar stablecoins has grown into a systemic layer of the crypto financial system, with basket-margin arrangements at centralized exchanges, index perpetuals on decentralized venues, and proof-of-reserves frameworks at regulated issuers all depending on the same object: a joint de-peg correlation surface that no market participant has ever actually observed---on a representative day, only three or four coins in a basket exhibit any peg deviation at all, and the correlation matrices published by clearinghouses, index providers, and regulators are filled in by smoothing over the gaps. The trouble is that the gaps are not statistical noise around a densely sampled surface; they are \emph{structural discontinuities}, and a smoother that cannot see the discontinuities systematically understates the joint tail risk it is meant to price.

Two edges of the same question follow. What is the sharpest description any correlation model can produce of the joint de-peg behavior that a stablecoin basket actually exhibits, as opposed to the smooth surface that basket-margin frameworks assume? And what is the intrinsic information content of a constraint variety on which peg-preserving equilibria fail on codimension-one strata organized by structural cause? The economic edge and the mathematical edge, it turns out, admit a common answer. That answer identifies a market-specific ceiling on what any smooth-surface method can achieve---a \emph{correlation-rigidity ceiling}\footnote{The notion of \emph{correlation rigidity} used here is structural---the local constancy of the co-movement structure within a regime, and its failure across a regime boundary---not the risk-management sense of ``rigid correlation matrix'' in \citet{Li2000} and the copula-parametric literature that follows it.}---and identifying its location is the object of this paper.

Three literatures speak to the question and none has resolved it. Stablecoin economics---the money-view foundations of \citet{GortonEtal2022, CatalinKirk2022}, empirical stability channels of \citet{LyonsViswanath2023, dAvernasEtal2024}, run dynamics of \citet{MaZengZhang2023}, noncustodial stability of \citet{Klages2022, KlagesMundtMinca2020}, general-equilibrium \citep{RoutledgeZetlin2022}, Terra analyses \citep{Uhlig2022, LiuLiSprenger2023}, and policy overlays \citep{FSB2023Stablecoin, BIS2023Stablecoin, PWG2021Stablecoin}---takes the peg-deviation process as given and studies its stability at the single-name level, not the cross-sectional correlation object that basket-margin rules price. Correlation baselines---Gaussian copula \citep{Li2000}, Vasicek single-factor \citep{Vasicek2002}, regime-switching \citep{Pelletier2006}, dynamic equicorrelation \citep{EngleKelly2012}, POET \citep{FanLiaoMincheva2013}, shrinkage \citep{LedoitWolf2004}, with the structural critique of \citet{CovalJurekStafford2009} and asymmetric-crisis evidence of \citet{AngChen2002}---impose a smooth correlation surface on a parametric family and do not bound how sharply correlation can be recovered when the constraint variety is naturally singular. Orthogonal to these, an implicit-atlas programme of virtual fundamental classes \citep{FukayaOno1999, LiTian1998, BehrendFantechi1997, Pardon2016, Pardon2019}, the wall-crossing formalism of \citet{KontsevichSoibelman2008, Bridgeland2007}, and the derived-critical-locus bridge of \citet{Joyce2015} produced the sharpest foundations for enumerative invariants on constraint varieties with codimension-one singular strata; to the author's knowledge, no prior work applies implicit-atlas VFCs to a financial-economic object.

Closing the gap requires an object that lives natively in both worlds. This paper represents the stablecoin issuer configuration as a point of a configuration space $\Mconf$ (\S\ref{sec:moduli}), whose peg-preserving locus is cut out by the peg constraint and whose strata are indexed by failure channel---backing-ratio, oracle, redemption, custody, bridge, regulatory, funding-rate, interest-rate. Across the boundary of a stratum the joint co-movement structure jumps by a rank-one term supported on the affected coins, and no smooth surface can reproduce that jump. The residue is stated on the co-movement \emph{covariance}, not on the correlation matrix, and the choice is forced rather than stylistic: a rank-one additive jump supported on a block of two or more coins necessarily adds a positive amount to its own diagonal, which a correlation matrix cannot absorb on either chamber side, since its diagonal is fixed at one. Carrying the crossing law on the covariance and passing to correlations by the normalising map removes the obstruction exactly, because that map \emph{assigns} the unit diagonal rather than computing it. Within a chamber of the peg-preserving parameter space, co-movement is locally constant and insensitive to smooth deformation of issuance rules; across a wall it jumps, and the jump has a closed form we state rather than estimate. Read economically, the three steps buy a single testable implication: the jump is \emph{supported on a small block of coins}, with a sign fixed by the wall taxonomy before the event window is observed, and it survives normalisation as a signed three-tier footprint on the correlation matrix rather than as a rank-one object. We call this footprint the \emph{block signature} of the wall---the coin pairs whose co-movement moves, and the sign in which it moves. This is the sense in which we use the word \emph{contagion}: not a diffuse rise in co-movement across the whole market, but a structured rise confined to the coins a failure channel touches, which is exactly what makes it identifiable. The block signature is what a supervisor reads off an acute-window correlation submatrix, and it is what every smooth-surface method is structurally unable to place there. That is what makes the diagnostic falsifiable rather than fitted.

The paper's diagnostic---basket risk is a staircase, not a slope---rests on three parallel contributions.

\textit{First, a structural representation of basket risk with a wall-stratified parameter space.} We represent the joint peg-preserving equilibrium as a configuration space---a parameter space whose geometry is inherited from the constraint that every issuer keep its peg---so that the stress boundaries unique to markets with structurally distinct failure triggers appear as strata rather than as features to be estimated. The construction reads directly off any public stablecoin panel and turns a two-dimensional smoothing exercise into a stratified recovery problem.

\textit{Second, a rigidity-and-crossing framework whose consequences are made explicit.} Three properties are posited and their consequences proved. The properties: co-movement constant within a failure regime up to a positive scale; a jump across a regime boundary that is a signed sum of rank-one terms; and support of that jump on the coins the boundary's channel touches. The consequences: correlation rigidity on chambers, with the scale cancelled exactly, so that rigidity survives normalisation without a further assumption; zero diagonal for the induced correlation jump; a signed three-tier footprint on the correlation image, replacing the rank-one structure that the normalisation destroys; and the gap $\pm\tfrac12\Deltait_R$ separating the true correlation from any locally continuous surface. Together these are what make basket risk a staircase rather than a slope---the question is not what the surface does between boundaries, because on a chamber it does nothing---and the passage to correlations tells us exactly what part of the block signature an estimator could even in principle target.

\textit{Third, a pair-level test of the block prediction.} Because the normalisation destroys rank-one structure and leaves a signed footprint, the empirical object is a gradient across pair layers: pairs of coins both inside the wall-affected block move up, pairs with one coin inside move the other way, and pairs with neither do not move at all. The block-versus-rest contrast differences out the crisis-wide co-movement that raises every layer at once. On a nine-event 2020--2025 panel built entirely from public data, the USDC--SVB custody event delivers the predicted block effect, with a block-versus-rest contrast of \GUSDCSVBDID\ and a coin-level permutation $p = \GUSDCSVBPCOIN$; the pair $(\text{DAI},\text{FRAX})$, which contains neither the primary victim nor any coin with direct custody exposure, jumps by \GUSDCSVBPURE---the largest of any pair---so the gradient is a block effect rather than a victim effect. The result survives a full $3\times3$ window grid. The test applies to the three events that carry a multi-coin block. The construction additionally identifies which policy instruments the block signature would key.

The remainder of the paper proceeds as follows. Section~\ref{sec:moduli} defines the stablecoin state, the configuration space and its peg-preserving locus, and states which properties of that space the paper assumes and which it derives. Section~\ref{sec:rigidity} defines co-movement as a second moment, takes the correlation matrix as its normalised image, and proves that rigidity of the covariance yields rigidity of the correlation matrix. Section~\ref{sec:wall} states the crossing law, derives its consequences for correlations, maps the eight stress walls, and contrasts them with the Gaussian copula. Section~\ref{sec:empirical} builds and runs the pair-level gradient test on the nine-event panel. Section~\ref{sec:conclusion} concludes; robustness checks are in the Appendix.
